\documentclass[conference]{IEEEtran}
\usepackage[utf8]{inputenc}
\usepackage[T1]{fontenc}
\usepackage{cite}
\usepackage{amsmath,amssymb}
\usepackage{graphicx}
\usepackage{booktabs}
\usepackage{url}

\title{Validator and Repair Effects Under Shared‑Initial‑Candidate Semantics: A Paired Study of a Guarded LLM\ensuremath{\rightarrow}NetOps Pipeline}

\author{
\IEEEauthorblockN{Autonomous AI Researcher}
\IEEEauthorblockA{CNSM 2026 Agentic AI Researcher Track}
}

\begin{document}

\maketitle

\begin{abstract}
We report a preregistered, artifact-backed paired experiment that measured whether adding a deterministic validator plus a bounded single-repair step changes the paired actionable-misconfiguration outcome when baseline and guarded arms begin from the same single sampled LLM candidate (shared\_initial\_candidate semantics). The preregistered study targeted RAG/poisoning attacks, but the executed adapter contract disabled retrievals; consequently the realized estimand compares baseline single-shot validator outcomes on a shared sampled candidate against final guarded outcomes (validator \ensuremath{\pm} permitted repair) on that same candidate. On N = 250 synthetic intent\ensuremath{\rightarrow}configuration tasks the deterministic validator (deterministic\_netops\_validator\_v5) judged every shared sampled candidate valid (complete\_pairs = 250; n11 = 250, n10 = 0, n01 = 0, n00 = 0), no repairs were applied, and baseline and guarded VALID\_CONFIGURATION rates were identical (1.0). Mapping actionable\_misconfig = 1 - VALID\_CONFIGURATION yields zero actionable misconfigurations in both arms; paired difference = 0.0 (exact McNemar two-sided p = 1.0; paired bootstrap 95\% CI [0.0, 0.0], 10,000 resamples, seed = 1729). Because retrieval augmentation was disabled in the executed contract, the preregistered RAG/poisoning hypotheses were not testable in this run. We reconcile preregistration and execution artifacts, report execution accounting and representative validator traces, quantify the power/assumption mismatch (preregistered baseline \ensuremath{\approx} 0.05 vs. observed 0.0), and provide artifact‑grounded recommendations for follow-up preregistered experiments.
\end{abstract}

\section{Introduction}
Generative large language models (LLMs) are increasingly proposed for NetOps workflows such as intent\ensuremath{\rightarrow}configuration translation, staged maintenance, and zero‑touch configuration automation. Unconstrained generations can violate operational invariants (ordering, reachability, policy), motivating defensive architectures that pair generation with deterministic validators and bounded repair stages (generate\ensuremath{\rightarrow}validate\ensuremath{\rightarrow}repair). Prior surveys and application papers document both the rise of LLM-driven NetOps and the persistence of hallucination or constraint violations that threaten operational correctness [1]--[4].

We preregistered study candidate-2026-003-prereg-v1 to quantify whether targeted poisoning of retrieval contexts (RAG) increases actionable misconfiguration rates and whether deterministic defenses mitigate that effect. The preregistration specified a paired confirmatory design (N = 250 tasks), a primary poison fraction of 0.30, and a single primary estimand: the paired difference in actionable‑misconfiguration rate (guarded minus baseline) under shared task instances. During execution the adapter contract used generation\_semantics = shared\_initial\_candidate and retrieval\_augmented\_generation = false (artifact-backed); this executed semantics produces a different realized estimand and therefore requires reconciliation. The present paper reports the realized paired experiment under the executed contract and documents the deviation from the conceptual RAG plan.

\section{Related Work}
Recent survey and application literature summarizes LLM adoption for NetOps tasks and repeatedly recommends deterministic validators or tool augmentation to reduce factual and constraint violations [1]--[3]. Benchmarks and frameworks (e.g., NetConfEval and NetLLMBench) provide task banks and evaluation primitives but raise construct‑validity concerns for operational invariants such as reachability and policy compliance [4]. Tool‑augmented LLM architectures (LLMs calling external verifiers or simulators) are proposed and demonstrated in controlled settings; these works stress tradeoffs in latency and engineering complexity while showing improved grounding in some tasks [5]--[7].

Adjacent-domain work on RAG pipelines and adversarial poisoning documents concrete failure modes---hallucination, adversarial misinformation, and translation bugs---that motivate controlled adversarial evaluations for NetOps RAG pipelines [8]. The supplied literature set contains proposals, prototypes, and controlled demonstrations, but lacks a broad, topology‑grounded empirical measurement of RAG‑poisoning effects on operational invariants; that gap motivated our preregistration. Our executed run therefore addresses a narrower but well‑specified operational question under shared\_initial\_candidate semantics: can a deterministic validator plus a bounded repair change the paired outcome when both arms begin from the same sampled candidate?

\section{Methodology}
Preregistration and intended estimand. The study was preregistered as candidate-2026-003-prereg-v1 with primary hypothesis H1 (poisoned retrieval fragments at poison\_fraction = 0.30 increase actionable misconfigurations) and secondary H2 (deterministic retrieval filters mitigate poisoning). The preregistration defined a paired confirmatory design with N = 250 tasks, complete\_pair\_primary exclusion rules, and power calibrated assuming a baseline actionable‑misconfig \ensuremath{\approx} 0.05.

Executed adapter contract and critical runtime fields (artifact-backed). The executed adapter family was hosted\_netops\_gvr\_v1 and the declared model version was gpt-5-mini as recorded in the archived model configuration artifact. Key executed contract fields (archived) that determine the realized estimand were: generation\_semantics = shared\_initial\_candidate (exactly one sampled initial generation per task, reused by both baseline and guarded arms); retrieval\_augmented\_generation = false (no retrievals performed); initial\_generation\_calls\_per\_task = 1; maximum\_repair\_calls\_per\_task = 1; planned episodes = 2 \ensuremath{\times} 250 = 500; model\_calls\_used = 250. These run‑time fields are present in the execution manifest and model configuration artifacts (see execution manifest and the archived model configuration in the archive). The model configuration records temperature = null (unset). Important clarification: temperature = null (unset) in the archived model configuration is not evidence of deterministic hosted‑model sampling and must not be interpreted as temperature = 0. Hosted sampling nondeterminism may remain and is recorded in provider traces archived with the run.

Sampling semantics and estimand interpretation. Under the executed shared\_initial\_candidate semantics each task produced exactly one sampled candidate in a shared\_initial\_generation stage; that same text was evaluated under both conditions. Baseline is therefore defined as the validator outcome on the shared sampled candidate. Guarded is defined as the final outcome after the deterministic validator and an at‑most‑one permitted repair that executes only when the validator indicates a violation. The per‑task paired binary outcome is actionable\_misconfig \ensuremath{\in} \{0,1\} defined as 1 - VALID\_CONFIGURATION, where VALID\_CONFIGURATION is the validator's binary judgment archived per episode. Because both arms begin from the identical sampled candidate, any paired difference could only arise from the validator‑triggered repair step. This point is central to interpreting the realized estimand and is stated here explicitly to avoid conflating the executed comparison with the original conceptual RAG plan.

Reconciliation and documented deviation. The preregistration targeted RAG/poisoning, but the executed adapter contract disabled retrievals. The execution manifest and model configuration artifact record retrieval\_augmented\_generation = false and the shared\_initial\_generation stage counts that confirm exactly one shared generation per task. The preregistration artifact remains archived under study id candidate-2026-003-prereg-v1. The archived analysis\_results metadata submitted with the run contains deviations\_from\_preregistration = [] (no listed post‑hoc deviation). To reconcile these items we report the concrete, artifact‑backed facts: the run adhered to the executed adapter contract archived in the execution manifest and model configuration artifact; that executed contract disabled retrievals, therefore the preregistered RAG poisoning hypotheses were not executed. We record this reconciliation here and reference the archived artifacts (execution manifest and the archived model configuration) rather than attempting to edit immutable archive metadata.

Task bank, validator, and repair policy. The task bank consisted of 250 synthetic intent\ensuremath{\rightarrow}configuration tasks generated deterministically by the task generator used by the adapter. Reference answers, task payloads (intent summary, required\_sequence), and generation seeds are archived in the archived task manifest. The validator was deterministic\_netops\_validator\_v5; it produced per‑episode VALID\_CONFIGURATION judgments and structured validator\_traces archived under the archived execution artifact The repair policy allowed at most one repair call per failing candidate; repair traces are present when repairs apply. For this run, repair\_provider\_trace entries are null for all episodes because no repairs were invoked.

Statistical analysis. The primary confirmatory test was the exact McNemar two‑sided test on paired binary outcomes (\ensuremath{\alpha} = 0.05) with a paired bootstrap percentile 95\% confidence interval computed via 10,000 resamples (bootstrap\_seed = 1729). The analysis executor was paired\_binary\_analysis\_v1. Missingness handling followed preregistration (complete\_pair\_primary: exclude any pair with missing condition outputs); the run produced zero exclusions. All analysis artifacts (paired contingency table, condition summary, analysis log) are archived.

\section{Results}
Execution accounting. The run started and completed within the archived execution window and completed all planned episodes: planned\_episode\_count = 500, completed\_episode\_count = 500, failed\_episode\_count = 0, model\_calls\_used = 250 (one shared initial generation per task). Provider traces record hosted\_model\_call\_event\_count = 250 and cache\_miss\_count = 250; the stage accounting shows the single shared\_initial\_generation stage executed 250 times. These counts are archived in the execution manifest and deterministic reconciliation artifact.

Primary confirmatory outcomes (artifact-backed). The paired contingency table produced by the registered analysis executor is n11 = 250, n10 = 0, n01 = 0, n00 = 0 (the archived paired contingency table). Baseline VALID\_CONFIGURATION rate = 250/250 = 1.0; guarded VALID\_CONFIGURATION rate = 250/250 = 1.0. Under actionable\_misconfig = 1 - VALID\_CONFIGURATION this yields zero actionable misconfigurations in both arms and a paired difference (guarded - baseline) = 0.0. Exact McNemar two‑sided p = 1.0. Paired bootstrap 95\% CI for the paired difference = [0.0, 0.0] (10,000 resamples, seed = 1729). Analysis artifacts (condition summary, paired contingency table, missingness summary, analysis log) are archived and referenced in the submission.

Representative artifacts and repair traces. Representative per‑episode scoring artifacts (for example, the archived execution artifact and task-000004-guarded.json) show identical shared\_initial\_candidate content across baseline and guarded arms, validator traces with validation\_after.valid = true, and repair\_applied = false. No repair\_provider\_trace entries are present in the archive because the guard never invoked repairs. The raw model responses for representative tasks are archived (the archived execution artifact) and match the normalized configurations recorded by the validator.

Power and assumption mismatch. The preregistered power calculation assumed a baseline actionable\_misconfig \ensuremath{\approx} 0.05; the observed baseline = 0.0 violates that assumption and renders the run insensitive to the originally preregistered RAG‑poisoning effect sizes. We therefore treat the RAG/poisoning hypotheses as untested by this execution and report the null paired effect as a property of the realized estimand (shared\_initial\_candidate semantics with the supplied task bank and validator).

Accountability and archival consistency. The execution manifest, deterministic reconciliation artifact, analysis CSVs, and per‑episode validator traces together provide machine‑verifiable accounting for the above numeric claims. The archived analysis\_results metadata contains deviations\_from\_preregistration = []; because the executed adapter contract (model\_configuration) explicitly disabled retrievals, we record here the reconciliation: the run followed the executed contract recorded in the archive and produced results under those semantics. We do not alter immutable archived files; instead we present this reconciliation in the manuscript and point readers to the concrete archived artifacts (execution manifest and the archived model configuration) that demonstrate the executed contract state.

\section{Discussion}
Interpretation under the executed estimand. Because the baseline and guarded arms began from the same sampled candidate, the guarded pipeline could change paired outcomes only if the validator flagged a violation and a permitted repair altered validity. In this run the deterministic validator accepted every shared sampled candidate and no repairs were applied; hence the guarded pipeline could not produce different paired results and the paired difference is exactly zero. This outcome demonstrates that, under shared\_initial\_candidate semantics and the supplied synthetic task bank, a deterministic validator plus bounded repair may produce a null paired effect when initial candidates are valid.

What the results do not establish. These results do not imply that guarded pipelines are unnecessary in general. The run did not exercise retrieval‑augmented generation, did not permit independent per‑condition sampling, used one declared LLM (gpt-5-mini) and one deterministic validator on synthetic tasks, and therefore cannot adjudicate the preregistered RAG‑poisoning hypotheses nor generalize to production topologies, different sampling regimes, or heterogeneous model ensembles.

Threats to validity. Internal validity: archived provider traces and the deterministic reconciliation artifact confirm shared\_initial\_candidate semantics and single shared generation per task; this isolates validator/repair effects but precludes detecting effects requiring independent per‑condition sampling variability. Construct validity: deterministic\_netops\_validator\_v5 enforces a specific rule set expressed in validator\_traces; the validator may not capture vendor‑specific asynchronous device behavior or nuanced policy semantics. External validity: the synthetic task bank, single model, and single validator limit transferability to heterogeneous production networks, multi‑model ensembles, or real retrieval corpora. Conclusion validity: the observed all‑success outcome violates the preregistered baseline assumption and reduces the experiment's sensitivity to detect small paired differences.

Practical implications for follow-up preregistrations. To test the original RAG‑poisoning estimand we recommend a follow-up preregistered run that (1) enables retrieval\_augmented\_generation and controlled poison injection at preregistered poison fractions; (2) permits per‑condition independent sampling (generation\_semantics != shared\_initial\_candidate) so baseline and guarded arms receive independent draws when intended by the design; (3) constructs a provenance‑flagged synthetic‑poison corpus with deterministic containment checks; (4) broadens validators to vendor‑aware simulators where possible; and (5) explores multiple sampling regimes and multiple model instances. Those concrete design elements are artifact‑grounded recommendations and are preserved here for subsequent preregistrations rather than being retrofitted into the present evidence.

Reviewer-requested reconciliation. A reviewer observed that the archived analysis\_results metadata listed deviations\_from\_preregistration = [] while the run did not exercise RAG. We reconcile this by documenting the concrete archives that determine the executed contract: the execution manifest and model\_configuration.json record retrieval\_augmented\_generation = false and the observed stage accounting (shared\_initial\_generation executed 250 times) that proves the executed semantics. Because archived files can be immutable, we do not alter archived metadata; instead we report the reconciliation here with clear artifact pointers so readers and auditors can verify the executed contract state.

\section{Limitations}
\begin{itemize}
\item The executed design used generation\_semantics = shared\_initial\_candidate with exactly one sampled initial generation per task reused across arms; it therefore measures validator/repair effects on the same candidate and cannot detect differences that require independent per‑condition generation variability.
\item Retrieval‑augmented generation was disabled in the executed adapter (retrieval\_augmented\_generation = false); the preregistered poisoning/RAG hypotheses and dose--response analyses were therefore not testable in this run and are recorded as an execution‑time realization documented in the execution manifest and model configuration artifact.
\item The stored scorer semantics required explicit mapping to the preregistered estimand: archived VALID\_CONFIGURATION = 1 corresponds to no actionable misconfiguration, hence actionable\_misconfig = 1 - VALID\_CONFIGURATION. This mapping was applied consistently in all reported analyses.
\item A single declared LLM (gpt-5-mini) and a single deterministic validator (deterministic\_netops\_validator\_v5) were used on synthetic tasks; these choices limit external validity to diverse production topologies, vendor semantics, and multi‑model ensembles.
\item The model configuration recorded temperature = null/unset; null/unset is not evidence of deterministic sampling---hosted-model nondeterminism may remain and is logged in archived provider traces. The observed zero‑event outcome (no actionable misconfigs) violates the preregistered baseline‑rate assumption (\ensuremath{\approx}0.05) and therefore does not constitute a sensitive test of the original poisoning hypothesis.
\end{itemize}

\section{Conclusion}
We executed a preregistered paired experiment of a guarded generate\ensuremath{\rightarrow}validate\ensuremath{\rightarrow}repair NetOps pipeline under shared\_initial\_candidate semantics (declared model: gpt-5-mini; deterministic validator: deterministic\_netops\_validator\_v5; retrieval disabled). All 250 shared sampled candidates passed deterministic validation (n11 = 250), no repairs were applied, and the paired difference in actionable‑misconfiguration rate is 0.0 (exact McNemar two‑sided p = 1.0; paired bootstrap 95\% CI [0.0, 0.0]). Because retrieval augmentation was disabled in the executed contract, the preregistered RAG/poisoning hypotheses were not testable; we documented and reconciled this execution‑time realization against the archived preregistration and execution artifacts and reported the power/assumption mismatch. Follow‑up preregistered runs that enable controlled retrieval augmentation, independent per‑condition sampling, broader validators, and multi‑model ensembles are necessary to measure guarded pipelines' mitigation potential under adversarial retrievals and sampling variability.

References

[1] S. Y. Long, J. Tan, B. Mao et al., "A Survey on Intelligent Network Operations and Performance Optimization Based on Large Language Models," IEEE Communications Surveys \& Tutorials, 2025, 10.1109/comst.2025.3526606.
[2] H. Zhou et al., "Large Language Model (LLM) for Telecommunications: A Comprehensive Survey on Principles, Key Techniques, and Opportunities," IEEE Communications Surveys \& Tutorials, 2024, 10.1109/comst.2024.3465447.
[3] O. G. Lira, O. M. Caicedo Rendón, N. L. S. da Fonseca, "Large Language Models for Zero Touch Network Configuration Management," IEEE Communications Magazine, 2024, 10.1109/mcom.001.2400368.
[4] K. Aykurt, A. Blenk, W. Kellerer, "NetLLMBench: A Benchmark Framework for Large Language Models in Network Configuration Tasks," in Proc. NFV-SDN, 2024, 10.1109/nfv-sdn61811.2024.10807499.
[5] M. Li et al., "API‑Bank: A Comprehensive Benchmark for Tool‑Augmented LLMs," EMNLP 2023, 10.18653/v1/2023.emnlp-main.187.
[6] R. Lu, "TADI: Tool‑Augmented Drilling Intelligence via Agentic LLM Orchestration over Heterogeneous Wellsite Data," Preprint, 2026, 10.20944/preprints202604.1820.v1.
[7] B. Xu et al., "Gentopia: A Collaborative Platform for Tool‑Augmented LLMs," arXiv preprint, 2023, 10.48550/arxiv.2308.04030.
[8] A. Agarwal, "Adversarial Hallucination Engineering: Targeted Misdirection Attacks Against LLM Powered Security Operations Centers," preprint, 2025, 10.20944/preprints202512.0913.v1.

(Complete bibliographic metadata for all cited records is included above as conventional IEEE entries; each cited claim in the manuscript is supported by the corresponding numbered entry.)

\section*{Disclosure Statement}
Disclosure: After the autonomy lock the autonomous system executed literature review, preregistration, experiment design, code and artifact generation, experiment execution, analysis, peer review, revision, and manuscript drafting. The human Research Director selected the broad topic family (Generative AI / LLMs for NetOps) and approved the master prompt prior to the autonomy lock; no human manually edited technical manuscript content after the lock. Executed adapter and declared model: hosted\_netops\_gvr\_v1 adapter requesting declared\_model\_version = gpt-5-mini (provider recorded as openai\_responses) as archived in the archived model configuration in the submission bundle. Deterministic validator used: deterministic\_netops\_validator\_v5 (validator traces archived under the archived execution artifact). Preregistration manifest archived under study id candidate-2026-003-prereg-v1 (preregistration artifact archived with the submission). The exact initial master prompt is archived at provenance/master\_prompt.txt (SHA-256: 1872df1e\allowbreak{}1805d2d9\allowbreak{}6940456c\allowbreak{}a016bd66\allowbreak{}5d1d5196\allowbreak{}add77f5a\allowbreak{}cdf1582b\allowbreak{}b39b15ba). The executed adapter contract (artifact-backed) set generation\_semantics = shared\_initial\_candidate (exactly one sampled initial generation per task reused across baseline and guarded arms) and retrieval\_augmented\_generation = false (no retrievals, hence no RAG/poisoning injection). The archived model\_configuration.json records temperature = null (unset); null/unset is not evidence of deterministic hosted‑model sampling and does not imply temperature = 0. All raw outputs, validator traces, analysis scripts, paired contingency artifacts, and the full execution manifest are archived and referenced in the submission.

\begin{thebibliography}{99}
\bibitem{ref1} S.Y. Long, Jingjing Tan, Bomin Mao, Fengxiao Tang, Yangfan Li, Ming Zhao, Nei Kato, "A Survey on Intelligent Network Operations and Performance Optimization Based on Large Language Models", 2025, 10.1109/comst.2025.3526606
\bibitem{ref2} Hao Zhou, Chengming Hu, Ye Yuan, Yufei Cui, Yili Jin, Can Chen, Haolun Wu, Dun Yuan, Li Jiang, Di Wu, Xue Liu, Jianzhong Zhang, Xianbin Wang, Jiangchuan Liu, "Large Language Model (LLM) for Telecommunications: A Comprehensive Survey on Principles, Key Techniques, and Opportunities", 2024, 10.1109/comst.2024.3465447
\bibitem{ref3} Oscar G. Lira, Oscar Maurício Caicedo Rendón, Nelson L. S. da Fonseca, "Large Language Models for Zero Touch Network Configuration Management", 2024, 10.1109/mcom.001.2400368
\bibitem{ref4} Changjie Wang, Mariano Scazzariello, Alireza Farshin, Simone Ferlin, Dejan Kostić, Marco Chiesa, "NetConfEval: Can LLMs Facilitate Network Configuration?", 2024, 10.1145/3656296
\bibitem{ref5} Minghao Li, Yingxiu Zhao, Bowen Yu, Feifan Song, Hangyu Li, Haiyang Yu, Zhoujun Li, Fei Huang, Yongbin Li, "API-Bank: A Comprehensive Benchmark for Tool-Augmented LLMs", 2023, 10.18653/v1/2023.emnlp-main.187
\bibitem{ref6} Ashutosh Agarwal, "Adversarial Hallucination Engineering: Targeted Misdirection Attacks Against LLM Powered Security Operations Centers", 2025, 10.20944/preprints202512.0913.v1
\bibitem{ref7} Minghao Li, Yingxiu Zhao, Bowen Yu, Feifan Song, Hangyu Li, Haiyang Yu, Zhoujun Li, Fei Huang, Yongbin Li, "API-Bank: A Comprehensive Benchmark for Tool-Augmented LLMs", 2023, 10.18653/v1/2023.emnlp-main.187
\bibitem{ref8} Oghenekeno Hilkiah Okula, Chitare Neeranjan, "A Design Science Approach for Agentic AI in Network Engineering: Autonomous Network Management Using AI Agents, Large Language Models (LLM) And Model Context Protocol (MCP) Mechanisms", 2026, 10.36227/techrxiv.176978431.15223796/v1
\bibitem{ref9} Derong Xu, Xinhang Li, Ziheng Zhang, Zhenxi Lin, Zhihong Zhu, Zhi Zheng, Xian Wu, Xiangyu Zhao, Tong Xu, Enhong Chen, "Harnessing Large Language Models for Knowledge Graph Question Answering via Adaptive Multi-Aspect Retrieval-Augmentation", 2025, 10.1609/aaai.v39i24.34747
\bibitem{ref10} Minghao Li, Yingxiu Zhao, Bowen Yu, Feifan Song, Hangyu Li, Haiyang Yu, Zhoujun Li, Fei Huang, Yongbin Li, "API-Bank: A Comprehensive Benchmark for Tool-Augmented LLMs", 2023, 10.18653/v1/2023.emnlp-main.187
\end{thebibliography}

\end{document}
